%% ===========================================================================
%%  Multimedia Tools and Applications (Springer Nature) -- sn-jnl.cls
%%  Evaluality: benchmark of no-reference quality scoring methods for the
%%  visual-quality selection criterion in multi-video summarization.
%%
%%  Build:  pdflatex main && bibtex main && pdflatex main && pdflatex main
%% ===========================================================================
%% MTAP requires sequential numbered citations in square brackets; the
%% reference list keeps the Springer Nature Basic format (sn-basic.bst).
\documentclass[sn-basic,Numbered]{sn-jnl}

\usepackage{graphicx}
\usepackage{multirow}
\usepackage{amsmath,amssymb,amsfonts}
\usepackage{booktabs}
\usepackage[table,dvipsnames]{xcolor}
\usepackage{manyfoot}      % required by sn-jnl.cls for \DeclareNewFootnote
\usepackage{hyperref}
\usepackage{tikz}
\usepackage{pgfplots}
\pgfplotsset{compat=1.17}
\usetikzlibrary{shapes.geometric, arrows.meta, positioning, calc, fit,
                backgrounds, patterns}
\usepackage{enumitem}
\usepackage{array}
\usepackage{threeparttable}
\usepackage{placeins}
\newcolumntype{L}[1]{>{\raggedright\arraybackslash}p{#1}}

%% ---------------------------------------------------------------------------
%%  VISUAL SYSTEM
%%  One palette for the whole paper: categorical hues encode *method family*
%%  and a divergent red-to-green ramp encodes *performance*. Both ramps are
%%  used consistently in every table and every figure.
%% ---------------------------------------------------------------------------

% Structural colours
\definecolor{headerblue}{HTML}{1B4F72}
\definecolor{rowlight}{HTML}{EDF3F8}
\definecolor{rowwhite}{HTML}{FFFFFF}
\definecolor{rulegray}{HTML}{9AA5B1}

% Performance ramp (red = poor, green = strong) -- used for cell shading
\definecolor{perf5}{HTML}{1E8449}   % strong          |SRCC| >= 0.70
\definecolor{perf4}{HTML}{7DCEA0}   % moderate        0.50 -- 0.70
\definecolor{perf3}{HTML}{F9E79F}   % weak            0.30 -- 0.50
\definecolor{perf2}{HTML}{F5B041}   % very weak       0.15 -- 0.30
\definecolor{perf1}{HTML}{E74C3C}   % negligible      < 0.15

% Accents
\definecolor{accentgreen}{HTML}{1E8449}
\definecolor{accentred}{HTML}{C0392B}
\definecolor{accentorange}{HTML}{E67E22}
\definecolor{accentblue}{HTML}{2471A3}

% Phase colours (methodological figures)
\definecolor{phase1color}{HTML}{2980B9}
\definecolor{phase2color}{HTML}{1E8449}
\definecolor{phase3color}{HTML}{7D3C98}
\definecolor{phase4color}{HTML}{E67E22}

% Categorical colours -- one per evaluation paradigm
\definecolor{catlmm}{HTML}{1B4F72}      % NR-IQA, LMM-based (reference model)
\definecolor{catclip}{HTML}{1E8449}     % NR-IQA, CLIP-based
\definecolor{catnss}{HTML}{C0392B}      % NR-IQA, NSS-based
\definecolor{catdeep}{HTML}{7D3C98}     % deep features, no BoVW
\definecolor{catitu}{HTML}{E67E22}      % ITU-T complexity
\definecolor{cathand}{HTML}{6C757D}     % handcrafted baseline
\definecolor{catbovw}{HTML}{C2185B}     % SIFT + BoVW (the deployed scoring method)
\definecolor{catcnnbovw}{HTML}{17A2B8}  % CNN + BoVW (ablation)
\definecolor{catcnn}{HTML}{0E6655}      % NR-IQA, CNN regression
\definecolor{cattrans}{HTML}{5B2C6F}    % NR-IQA, transformer

%% Performance-shaded table cells. \pfA is the strongest band, \pfE the weakest.
\newcommand{\pfA}[1]{\cellcolor{perf5}\textcolor{white}{\textbf{#1}}}
\newcommand{\pfB}[1]{\cellcolor{perf4}#1}
\newcommand{\pfC}[1]{\cellcolor{perf3}#1}
\newcommand{\pfD}[1]{\cellcolor{perf2}#1}
\newcommand{\pfE}[1]{\cellcolor{perf1}\textcolor{white}{#1}}
\newcommand{\pfNA}{\cellcolor{gray!12}\textcolor{gray!70}{n/a}}

%% Column head on the dark band
\newcommand{\hd}[1]{\textcolor{white}{\textbf{#1}}}

%% Score polarity markers (see Section 3.4)
\newcommand{\polup}{\ensuremath{\uparrow}}
\newcommand{\poldown}{\ensuremath{\downarrow}}

\begin{document}

\tikzset{
  phasebox/.style 2 args={rectangle, rounded corners=4pt, minimum width=2.9cm,
    minimum height=2.5cm, text width=2.7cm, align=center, font=\sffamily\small,
    draw=#1!80, fill=#1!10, line width=0.9pt, #2},
  iolabel/.style={font=\sffamily\scriptsize, text width=2.9cm, align=center},
  flowarr/.style={-{Stealth[length=6pt, width=5pt]}, thick, rulegray},
}


\title[Supplementary material]{Electronic Supplementary Material for
``A Component-Level Audit of No-Reference Quality Scoring in a Multi-Video
Summarization Pipeline''}

\author*[1]{\fnm{Jader L.} \sur{de Souza Gon\c{c}alves}}\email{jaderlouis@proton.me}
\author[1]{\fnm{Rudinei} \sur{Goularte}}\email{rudinei@icmc.usp.br}
\affil[1]{\orgdiv{Instituto de Ci\^{e}ncias Matem\'{a}ticas e de Computa\c{c}\~{a}o},
  \orgname{Universidade de S\~{a}o Paulo (USP)},
  \orgaddress{\city{S\~{a}o Carlos}, \state{SP}, \country{Brazil}}}

\maketitle

\renewcommand{\thetable}{S\arabic{table}}
\renewcommand{\theHtable}{S\arabic{table}}
\setcounter{table}{0}

This supplement holds the detailed breakdowns behind results reported in the
article, relocated here to keep the article at a readable length. Nothing in it
is new: every figure quoted in the article is reproduced in full below, and the
article cites these tables as Table~S1, Table~S2, and so on. Values, polarity
arrows, labels and boldface make every colour cue redundant in grayscale.

%% --- S1: tab:dataset ---
\begin{table}[!ht]
\centering
\caption{Human-annotated sensitivity collection (15 clusters): Brazilian television news topics, with
number of source videos, segment slots, redundancy clusters and segment
duration range. \emph{Slots} counts cluster memberships; because two segments
belong to two clusters each, the 33 slots correspond to 31 distinct segments.}
\label{tab:dataset}
\small
\setlength{\tabcolsep}{6pt}
\renewcommand{\arraystretch}{1.12}
\begin{tabular}{@{}l c c c c l@{}}
\toprule
\rowcolor{headerblue}
\hd{Topic} & \hd{Videos} & \hd{Slots} & \hd{Clusters} &
\hd{Duration (s)} & \hd{Status} \\
\midrule
\rowcolor{rowlight} \texttt{bebe\_real}            & 3 & 5 & 2 & 8--41   & Included \\
\rowcolor{rowwhite} \texttt{brumadinho}            & 3 & 6 & 3 & 13--188 & Included \\
\rowcolor{rowlight} \texttt{fone\_canc.\_ruido}    & 3 & 6 & 3 & 16--30  & Included \\
\rowcolor{rowwhite} \texttt{isencao\_imposto}      & 3 & 6 & 3 & 25--30  & Included \\
\rowcolor{rowlight} \texttt{mounjaro}              & 3 & 4 & 1 & 8--25   & Included \\
\rowcolor{rowwhite} \texttt{ucrania}               & 3 & 6 & 3 & 25--29  & Included \\
\midrule
\rowcolor{gray!10} \texttt{esquema\_pcc}           & 3 & -- & -- & -- &
  \textcolor{accentred}{\textbf{Excluded}}$^{\dagger}$ \\
\midrule
\rowcolor{rowlight} \textbf{Total} & \textbf{18} & \textbf{33} & \textbf{15} &
  8--188 & --- \\
\bottomrule
\multicolumn{6}{@{}p{12.4cm}}{\scriptsize $^{\dagger}$This six-topic set excludes
\texttt{esquema\_pcc}: 98 SIFT descriptors $<K=300$. The operational block
below is a distinct 83-segment redundancy output with different boundaries and
composition; its six \texttt{esquema\_pcc} segments form three valid clusters.
The sets are not nested. Slots exceed distinct segments where a segment belongs
to two clusters.} \\
\end{tabular}
\end{table}

\FloatBarrier
\begin{center}
\textbf{Table S1 (continued). Operational 42-cluster set used in the primary
application analysis.}\par\smallskip
\scriptsize
\setlength{\tabcolsep}{5pt}
\begin{tabular}{@{}lrrr@{}}
\toprule
\rowcolor{headerblue}
\hd{Topic} & \hd{Clusters} & \hd{Slots} & \hd{Distinct segments} \\
\midrule
\texttt{acordo\_mariana} & 1 & 2 & 2 \\
\rowcolor{rowlight}\texttt{aumento\_cafe} & 1 & 4 & 4 \\
\texttt{bebe\_real} & 2 & 5 & 5 \\
\rowcolor{rowlight}\texttt{bolsonaro\_trump} & 3 & 6 & 6 \\
\texttt{brasil\_omc} & 1 & 2 & 2 \\
\rowcolor{rowlight}\texttt{brumadinho} & 3 & 6 & 5 \\
\texttt{catedral\_notre\_dame} & 2 & 5 & 4 \\
\rowcolor{rowlight}\texttt{cessar\_fogo\_ucrania} & 3 & 6 & 4 \\
\texttt{esquema\_pcc} & 3 & 6 & 6 \\
\rowcolor{rowlight}\texttt{fone\_cancelamento\_ruido} & 3 & 6 & 6 \\
\texttt{fraude\_inss} & 3 & 6 & 4 \\
\rowcolor{rowlight}\texttt{google\_huawei} & 3 & 6 & 4 \\
\texttt{isencao\_imposto} & 3 & 6 & 5 \\
\rowcolor{rowlight}\texttt{julgamento\_bolsonaro} & 3 & 6 & 5 \\
\texttt{mounjaro} & 1 & 4 & 4 \\
\rowcolor{rowlight}\texttt{recorde\_calor\_2024} & 2 & 5 & 5 \\
\texttt{rio\_celular} & 1 & 2 & 2 \\
\rowcolor{rowlight}\texttt{ucrania} & 3 & 6 & 6 \\
\texttt{ultraprocessados} & 1 & 4 & 4 \\
\midrule
\textbf{Total} & \textbf{42} & \textbf{93} & \textbf{83} \\
\bottomrule
\end{tabular}
\end{center}


%% --- S2: tab:degradation ---
\begin{table}[!ht]
\centering
\caption{Illustration of the sensitivity axis on a single broadcast frame and
four degradations, after applying each method's declared polarity. The pattern
this small case shows is the one Table~\ref{tab:sensitivity} then measures over
$25$ distortion types and $10{,}125$ images; nothing is concluded from this
table alone. Cells are shaded by outcome:
\colorbox{perf5}{\textcolor{white}{expected order}},
\colorbox{perf3}{partial}, \colorbox{perf1}{\textcolor{white}{violated}}.}
\label{tab:degradation}
\small
\setlength{\tabcolsep}{5.5pt}
\renewcommand{\arraystretch}{1.14}
\begin{tabular}{@{}l c r r r r r l@{}}
\toprule
\rowcolor{headerblue}
\hd{Method} & \hd{Pol.} & \hd{Original} & \hd{Blur} & \hd{Noise} &
\hd{Dark} & \hd{JPEG} & \hd{Outcome} \\
\midrule
\rowcolor{rowlight} QualiCLIP & \polup & 0.341 & 0.155 & 0.167 & 0.196 & 0.163
  & \pfA{as expected} \\
\rowcolor{rowwhite} LIQE      & \polup & 1.294 & 1.056 & 1.000 & 1.023 & 1.023
  & \pfA{as expected} \\
\rowcolor{rowlight} BRISQUE   & \poldown & 42.8 & 90.0 & 66.4 & 59.1 & 78.2
  & \pfA{as expected} \\
\rowcolor{rowwhite} NIQE      & \poldown & 4.53 & 11.99 & 15.79 & 6.18 & 9.47
  & \pfA{as expected} \\
\midrule
\rowcolor{catbovw!12} Keypoint density & \polup & 2598 & 104 & \textbf{4543} & 4
  & 2799 & \pfC{partial} \\
\midrule
\rowcolor{rowwhite} DINOv2 ViT-S/14 & \polup & 14.69 & 14.86 & 14.86 & 14.90
  & 15.14 & \pfE{violated} \\
\rowcolor{rowlight} ResNet-50 & \polup & 15.62 & \textbf{18.11} & 15.00 & 15.97
  & 15.28 & \pfE{violated} \\
\bottomrule
\multicolumn{8}{@{}l}{\scriptsize Values in bold are the ones that violate the
expected order.} \\
\end{tabular}
\end{table}

\FloatBarrier
\begin{center}
\textbf{Table S2 (continued). Exact parameters of the constructed broadcast
degradations.}\par\smallskip
\scriptsize
\setlength{\tabcolsep}{4pt}
\begin{tabular}{@{}lp{8.8cm}@{}}
\toprule
\rowcolor{headerblue}\hd{Type} & \hd{Levels 1--5} \\
\midrule
Gaussian blur & Kernel width $\{3,5,9,15,23\}$ pixels \\
\rowcolor{rowlight}White noise & Gaussian $\sigma=\{5,12,22,35,50\}$; deterministic seeds $1001$--$1005$ \\
Darkening & RGB multiplier $\{.80,.62,.46,.32,.20\}$ with clipping \\
\rowcolor{rowlight}JPEG & OpenCV quality $\{40,25,15,9,5\}$ \\
H.264 & libx264 CRF $\{30,36,41,45,50\}$, preset \texttt{veryfast} \\
\rowcolor{rowlight}Rescale & Downscale factors $\{.75,.55,.40,.28,.18\}$, then restore original size \\
Motion blur & Horizontal kernel width $\{5,9,15,23,33\}$ pixels \\
\rowcolor{rowlight}Desaturation & HSV saturation multiplier $\{.75,.55,.35,.18,.05\}$ \\
\bottomrule
\multicolumn{2}{@{}p{0.96\linewidth}}{\scriptsize Each segment contributes at
most 30 uniformly spaced frames. The seeded 20-segment simple random sample
represents 13 topics: one segment each from \texttt{aumento\_cafe},
\texttt{bolsonaro\_trump}, \texttt{cessar\_fogo\_ucrania},
\texttt{esquema\_pcc}, \texttt{isencao\_imposto}, \texttt{mounjaro} and
\texttt{rio\_celular}; two each from \texttt{bebe\_real},
\texttt{google\_huawei}, \texttt{julgamento\_bolsonaro},
\texttt{recorde\_calor\_2024} and \texttt{ucrania}; and three from
\texttt{fone\_cancelamento\_ruido}.} \\
\end{tabular}
\end{center}


%% --- S3: tab:sensitivity ---
\begin{table}[!ht]
\centering
\caption{Sensitivity to distortion level on KADID-10k. Direction-aligned SRCC is
computed within each of the 81 reference images across five levels and then
averaged by distortion type and across the 25 types. Intervals block-bootstrap
the 81 references ($10^4$ resamples, seed 42). The last column counts types whose
reference-averaged SRCC exceeds $0.50$.}
\label{tab:sensitivity}
\small
\setlength{\tabcolsep}{6pt}
\renewcommand{\arraystretch}{1.14}
\begin{tabular}{@{}l r l r@{}}
\toprule
\rowcolor{headerblue}
\hd{Method} & \hd{Mean $\rho^{\ast}$} & \hd{$95\%$ block CI} & \hd{Types $>0.5$} \\
\midrule
\rowcolor{rowlight} Q-Align   & \pfA{0.880} & $[0.874,0.886]$ & \pfA{23/25} \\
\rowcolor{rowwhite} QualiCLIP & \pfA{0.849} & $[0.840,0.857]$ & \pfA{22/25} \\
\rowcolor{rowlight} CLIP-IQA$+$ & \pfA{0.808} & $[0.795,0.821]$ & \pfA{22/25} \\
\rowcolor{rowwhite} LIQE      & \pfA{0.807} & $[0.797,0.817]$ & \pfA{22/25} \\
\rowcolor{rowlight} MUSIQ     & \pfA{0.733} & $[0.720,0.747]$ & \pfA{20/25} \\
\rowcolor{rowwhite} TOPIQ     & \pfA{0.721} & $[0.702,0.737]$ & \pfA{20/25} \\
\rowcolor{rowlight} HyperIQA  & \pfB{0.673} & $[0.653,0.692]$ & \pfA{19/25} \\
\rowcolor{rowwhite} DBCNN     & \pfB{0.668} & $[0.642,0.692]$ & \pfA{19/25} \\
\rowcolor{rowlight} ARNIQA$^{\dagger}$ & \pfB{0.668} & $[0.652,0.683]$ & \pfA{20/25} \\
\rowcolor{rowwhite} BRISQUE   & \pfB{0.657} & $[0.628,0.684]$ & \pfA{19/25} \\
\rowcolor{rowlight} CLIP-IQA  & \pfB{0.645} & $[0.612,0.676]$ & \pfA{20/25} \\
\rowcolor{rowwhite} TReS$^{\dagger}$ & \pfB{0.620} & $[0.598,0.640]$ & \pfA{18/25} \\
\rowcolor{rowlight} NIQE      & \pfB{0.607} & $[0.587,0.625]$ & \pfA{19/25} \\
\midrule
\rowcolor{rowwhite} SI              & \pfD{0.251} & $[0.219,0.282]$ & \pfD{13/25} \\
\rowcolor{catbovw!12} Keypoint density & \pfE{0.079} & $[0.033,0.128]$ & \pfE{\ \,7/25} \\
\rowcolor{catbovw!12} SIFT+BoVW       & \pfE{$-0.077$} & $[-0.124,-0.030]$ & \pfE{\ \,5/25} \\
\rowcolor{rowwhite} DINOv2          & \pfE{$-0.064$} & $[-0.136,0.009]$ & \pfE{\ \,1/25} \\
\rowcolor{rowlight} ResNet-50       & \pfE{$-0.297$} & $[-0.364,-0.225]$ & \pfE{\ \,0/25} \\
\bottomrule
\multicolumn{4}{@{}p{0.94\linewidth}}{\scriptsize $^{\dagger}$ARNIQA and TReS are
reported on this axis as coverage extensions but are not treated as complete-case
methods in the screening rule because the required CLIVE evaluations were not available.} \\
\end{tabular}
\end{table}


%% --- S4: tab:p1401 ---
\begin{table}[!ht]
\centering
\caption{P.1401 point report after the logistic mapping. PLCC measures linearity,
RMSE accuracy and the outlier ratio (OR; lower is better) consistency;
RMSE$^{*}$ is the companion error after discounting the $95\%$ confidence
interval of the subjective score. OR saturates on KonIQ-10k because its
per-image subjective intervals are narrow, but remains discriminative on
KADID-10k. Analytic $95\%$ intervals for PLCC, RMSE, RMSE$^{*}$ and OR are
included in the released machine-readable results. A dagger marks evaluation
on training data.}
\label{tab:p1401}
\scriptsize
\setlength{\tabcolsep}{2.5pt}
\renewcommand{\arraystretch}{1.10}
\begin{tabular*}{\textwidth}{@{\extracolsep{\fill}}l r r r r r r r r@{}}
\toprule
\rowcolor{headerblue}
\hd{Method} & \multicolumn{4}{c}{\cellcolor{headerblue}\hd{KonIQ-10k}} &
\multicolumn{4}{c}{\cellcolor{headerblue}\hd{KADID-10k}} \\
\rowcolor{headerblue}
& \hd{PLCC} & \hd{RMSE} & \hd{RMSE$^{*}$} & \hd{OR}
& \hd{PLCC} & \hd{RMSE} & \hd{RMSE$^{*}$} & \hd{OR} \\
\midrule
\rowcolor{rowlight} Q-Align   & .940 & 5.251 & 5.166 & .983 & .952 & .330 & .155 & .292 \\
\rowcolor{rowwhite} LIQE      & .915 & 6.234 & 6.149 & .983 & .720 & .751 & .526 & .703 \\
\rowcolor{rowlight} QualiCLIP & .771 & 9.837 & 9.752 & .990 & .673 & .800 & .572 & .725 \\
\rowcolor{rowwhite} TOPIQ$^{\dagger}$    & .978 & 3.198 & 3.124 & .959 & .546 & .907 & .671 & .769 \\
\rowcolor{rowlight} HyperIQA$^{\dagger}$ & .956 & 4.507 & 4.428 & .977 & .503 & .936 & .700 & .776 \\
\rowcolor{rowwhite} DBCNN$^{\dagger}$    & .935 & 5.483 & 5.399 & .982 & .499 & .939 & .700 & .799 \\
\rowcolor{rowlight} MUSIQ$^{\dagger}$    & .904 & 6.612 & 6.527 & .985 & .596 & .869 & .641 & .730 \\
\rowcolor{rowwhite} CLIP-IQA$+^{\dagger}$& .863 & 7.790 & 7.705 & .987 & .665 & .809 & .581 & .727 \\
\rowcolor{rowlight} CLIP-IQA  & .724 & 10.647 & 10.561 & .992 & .521 & .924 & .691 & .787 \\
\rowcolor{rowwhite} SI        & .561 & 12.782 & 12.696 & .993 & .328 & 1.023 & .783 & .821 \\
\rowcolor{rowlight} SIFT+BoVW (canon.) & .484 & 13.508 & 13.423 & .995 & .224 & 1.055 & .814 & .835 \\
\rowcolor{rowwhite} Keypoint density   & .484 & 13.502 & 13.417 & .994 & .227 & 1.055 & .814 & .836 \\
\rowcolor{rowlight} NIQE      & .393 & 14.194 & 14.108 & .993 & .439 & .973 & .734 & .807 \\
\rowcolor{rowwhite} BRISQUE   & .253 & 14.931 & 14.844 & .995 & .395 & .995 & .758 & .807 \\
\rowcolor{rowlight} ResNet-50 & .162 & 15.231 & 15.143 & .994 & .194 & 1.062 & .821 & .834 \\
\rowcolor{rowwhite} DINOv2    & .065 & 15.402 & 15.315 & .996 & .041 & 1.082 & .841 & .840 \\
\bottomrule
\end{tabular*}
\end{table}


%% --- S5: tab:cat ---
\begin{table}[!ht]
\centering
\caption{Mean direction-aligned SRCC by paradigm. The two blocks are separated
because the change between databases means different things in each. In the
upper block no member was fitted on either database, so $\Delta$ compares two
generalization results. In the lower block at least one member was fitted on
KonIQ-10k, marked $\dagger$; there the KonIQ-10k column mixes fit with
generalization and no $\Delta$ is meaningful.}
\label{tab:cat}
\footnotesize
\setlength{\tabcolsep}{4pt}
\renewcommand{\arraystretch}{1.14}
\begin{tabular}{@{}l c r r r l@{}}
\toprule
\rowcolor{headerblue}
\hd{Paradigm} & \hd{$N$} & \hd{KADID} & \hd{KonIQ} & \hd{$\Delta$} &
\hd{Best member} \\
\midrule
\multicolumn{6}{@{}l}{\itshape\textcolor{headerblue}{Neither database seen by any member}} \\
\rowcolor{rowlight} ITU-T complexity   & 1  & \pfD{0.219} & \pfB{0.513} & $+0.294$ & SI \\
\rowcolor{rowwhite} SIFT+BoVW$^{\ast}$ & 2  & \pfE{0.124} & \pfC{0.446} & $+0.322$ & canonical IDF \\
\rowcolor{rowlight} Handcrafted        & 1  & \pfE{0.123} & \pfC{0.447} & $+0.324$ & Keypoint density \\
\rowcolor{rowwhite} NR-IQA, NSS        & 2  & \pfC{0.355} & \pfC{0.302} & $-0.053$ & NIQE \\
\rowcolor{rowlight} CNN$+$BoVW         & 19 & \pfC{0.318} & \pfE{0.058} & $-0.260$ & VGG-19$+$BoVW \\
\rowcolor{rowwhite} Deep features      & 2  & \pfE{$-$0.086} & \pfE{$-$0.115} & $-0.029$ & DINOv2 ViT-S/14 \\
\midrule
\multicolumn{6}{@{}l}{\itshape\textcolor{headerblue}{At least one member fitted on KonIQ-10k}} \\
\rowcolor{rowlight} NR-IQA, LMM        & 1  & \pfA{0.951$^{\dagger}$} & \pfA{0.920$^{\dagger}$} & --- & Q-Align \\
\rowcolor{rowwhite} NR-IQA, CNN reg.   & 2  & \pfC{0.473} & \pfA{0.933$^{\dagger}$} & --- & HyperIQA \\
\rowcolor{rowlight} NR-IQA, transformer& 2  & \pfB{0.534} & \pfA{0.923$^{\dagger}$} & --- & TOPIQ \\
\rowcolor{rowwhite} NR-IQA, CLIP       & 4  & \pfB{0.636} & \pfA{0.774$^{\dagger}$} & --- & LIQE \\
\bottomrule
\multicolumn{6}{@{}p{0.92\linewidth}}{\scriptsize $^{\ast}$Mean over the
magnitudes of the two TF-IDF variants, which differ in orientation but not in
strength. Restricted to the two CLIP members fitted on neither database, that
paradigm gives $0.583$ on KADID-10k and $0.679$ on KonIQ-10k, a change of
$+0.096$.} \\
\end{tabular}
\end{table}


%% --- S6: tab:variants ---
\begin{table}[!ht]
\centering
\caption{The BoVW variants in the domain, on the agreement axis and on what they
deliver. \emph{Differs} counts the clusters in which the variant selects a
different segment from the deployed one; \emph{Agreement} is top-1 agreement with
the reference model over the 15 clusters; $r_{\mathrm{kp}}$ is the rank
correlation with keypoint density over the 33 segment slots; and the last three
columns are the summary the variant delivers, as in Table~\ref{tab:e2e}. The
comparison between summing and averaging defined in the main article is
therefore evaluated directly. The final column is the uncalibrated raw sum of
the DOVER technical and aesthetic branches, used only as a scale-sensitivity
diagnostic.}
\label{tab:variants}
\footnotesize
\setlength{\tabcolsep}{3.5pt}
\renewcommand{\arraystretch}{1.14}
\begin{tabular}{@{}l c c c r r r r@{}}
\toprule
\rowcolor{headerblue}
\hd{Variant} & \hd{Aggr.} & \hd{Differs} & \hd{Agreem.} &
\hd{$r_{\mathrm{kp}}$} & \hd{Dur.\ (s)} & \hd{Q-Align} & \hd{raw T+A} \\
\midrule
\rowcolor{catbovw!12} SIFT+BoVW (as implemented) & sum & --- & \pfD{9/15} & $+0.591$ & $634.0$ & \pfC{$61.1\%$} & \pfD{$52.1\%$} \\
\rowcolor{rowwhite} SIFT+BoVW (canonical IDF)  & sum & 1/15 & \pfC{10/15} & $+0.116$ & $801.5$ & \pfC{$62.2\%$} & \pfE{$47.6\%$} \\
\rowcolor{rowlight} BoVW / descriptor count    & mean & 2/15 & \pfC{10/15} & $-0.199$ & $810.9$ & \pfC{$59.0\%$} & \pfE{$48.2\%$} \\
\rowcolor{perf4!22} \textbf{BoVW / frame count} & mean & 1/15 & \pfC{10/15} & $-0.145$ & $\mathbf{459.3}$ & \pfB{$\mathbf{77.3\%}$} & \pfB{$\mathbf{71.5\%}$} \\
\midrule
\rowcolor{rowlight} Keypoint density           & mean & 5/15 & \pfE{\ \,7/15} & $+1.000$ & --- & --- & --- \\
\rowcolor{perf5!16} \textbf{LIQE} (for reference) & mean & 5/15 & \pfB{11/15} & $-0.381$ & $455.6$ & \pfA{$85.8\%$} & \pfA{$91.8\%$} \\
\bottomrule
\end{tabular}
\end{table}


%% --- S7: tab:agree ---
\begin{table}[!ht]
\centering
\caption{Descriptive in-domain agreement with the Q-Align reference model over
the 15 human-annotated redundancy clusters, raw and chance-corrected against
$\bar{p}_e = 0.472$. These clusters share topics and two pairs share a segment;
the table therefore makes no cluster-independence inference. Topic-block
intervals for the larger system-defined set are reported in the main article.
Roles follow Table~1: D = deployed, C = candidate and X = diagnostic.}
\label{tab:agree}
\small
\setlength{\tabcolsep}{6pt}
\renewcommand{\arraystretch}{1.14}
\begin{tabular}{@{}l c r r@{}}
\toprule
\rowcolor{headerblue}
\hd{Method} & \hd{Role} & \hd{Agreement} & \hd{$\kappa_{\mathrm{uniform}}$} \\
\midrule
\rowcolor{rowlight} MUSIQ     & C & \pfB{13/15 \ (86.7\%)} & \pfB{$+0.747$} \\
\rowcolor{rowwhite} DBCNN     & C & \pfB{12/15 \ (80.0\%)} & \pfB{$+0.621$} \\
\rowcolor{rowlight} TOPIQ     & C & \pfB{12/15 \ (80.0\%)} & \pfB{$+0.621$} \\
\rowcolor{rowwhite} LIQE      & C & \pfB{11/15 \ (73.3\%)} & \pfB{$+0.495$} \\
\rowcolor{rowlight} HyperIQA  & C & \pfB{11/15 \ (73.3\%)} & \pfB{$+0.495$} \\
\rowcolor{rowwhite} NIQE      & C & \pfC{10/15 \ (66.7\%)} & \pfC{$+0.368$} \\
\rowcolor{rowlight} QualiCLIP & C & \pfC{10/15 \ (66.7\%)} & \pfC{$+0.368$} \\
\rowcolor{rowwhite} BRISQUE   & C & \pfC{10/15 \ (66.7\%)} & \pfC{$+0.368$} \\
\rowcolor{rowlight} SIFT+BoVW & D & \pfD{\ \,9/15 \ (60.0\%)} & \pfD{$+0.242$} \\
\rowcolor{rowwhite} CLIP-IQA  & C & \pfD{\ \,8/15 \ (53.3\%)} & \pfD{$+0.116$} \\
\rowcolor{rowlight} CLIP-IQA$+$ & C & \pfD{\ \,8/15 \ (53.3\%)} & \pfD{$+0.116$} \\
\rowcolor{rowwhite} DINOv2    & X & \pfE{\ \,5/15 \ (33.3\%)} & \pfE{$-0.263$} \\
\bottomrule
\end{tabular}
\end{table}


%% --- S8: tab:divergence ---
\begin{table}[!ht]
\centering
\caption{The four redundancy clusters in which the method under study
diverges from the majority of the quality references. Each row is a candidate
segment; each column a scoring method. The top-1 choice of each method is
shaded---\colorbox{catbovw!22}{pink} for SIFT+BoVW,
\colorbox{perf5!22}{green} for the quality references. NIQE and BRISQUE are
lower-is-better. Scores are reproduced from the aggregated per-cluster results.}
\label{tab:divergence}
\scriptsize
\setlength{\tabcolsep}{1.5pt}
\renewcommand{\arraystretch}{1.12}
\newcommand{\pS}[1]{\cellcolor{catbovw!22}\textbf{#1}}
\newcommand{\pR}[1]{\cellcolor{perf5!22}\textbf{#1}}
\begin{tabular}{@{}l r r r r r r@{}}
\toprule
\rowcolor{headerblue}
\hd{Segment} & \hd{SIFT+BoVW} & \hd{LIQE} &
\hd{NIQE\,\poldown} & \hd{Q-Align} & \hd{QualiCLIP} & \hd{BRISQUE\,\poldown} \\
\midrule
\multicolumn{7}{@{}l}{\itshape\textcolor{headerblue}{\texttt{bebe\_real}~C1}} \\
Band \ 0:37 (10\,s)   & \pS{5448.0} & 1.49 & 6.55 & 4.02 & 0.28 & 50.97 \\
Record \ 0:46 (11\,s) & 2077.6 & \pR{1.59} & \pR{5.57} & \pR{4.17} & \pR{0.30} & \pR{41.44} \\
\midrule
\multicolumn{7}{@{}l}{\itshape\textcolor{headerblue}{\texttt{brumadinho}~C2}} \\
Band \ 0:00 (13\,s)    & 2888.3 & \pR{2.60} & \pR{5.04} & \pR{4.84} & \pR{0.36} & \pR{44.95} \\
Record \ 4:24 (188\,s) & \pS{5332.7} & 1.13 & 6.23 & 3.74 & 0.27 & 53.57 \\
\midrule
\multicolumn{7}{@{}l}{\itshape\textcolor{headerblue}{\texttt{isencao\_imposto}~C0}} \\
Band \ 2:46 (30\,s) & \pS{3732.8} & 4.23 & \pR{4.67} & 3.70 & \pR{0.62} & 34.65 \\
CNN \ 1:12 (25\,s)  & 3010.0 & \pR{4.47} & 4.74 & \pR{4.26} & 0.59 & \pR{29.84} \\
\midrule
\multicolumn{7}{@{}l}{\itshape\textcolor{headerblue}{\texttt{mounjaro}~C0}} \\
Band \ 0:00 (8\,s)       & 559.8  & 4.12 & 7.29 & \pR{4.38} & 0.56 & 63.19 \\
Dom.\,Esp.\ 2:06 (16\,s) & \pS{1945.1} & 3.95 & 5.87 & 3.90 & 0.57 & 36.55 \\
Dom.\,Esp.\ 3:07 (25\,s) & 1784.4 & 3.61 & 5.28 & 3.68 & 0.56 & 42.29 \\
Record \ 0:00 (14\,s)    & 815.2  & \pR{4.93} & \pR{4.76} & 4.26 & \pR{0.63} & \pR{35.45} \\
\bottomrule
\end{tabular}
\end{table}


%% --- S9: tab:pooling ---
\begin{table}[!ht]
\centering
\caption{Top-1 agreement with the reference model over the 15 clusters under
four temporal poolings. The worst-frame, tenth-percentile and
$\mathrm{mean}-\mathrm{sd}$ rules all penalise within-segment dispersion; the
mean ignores it.}
\label{tab:pooling}
\small
\setlength{\tabcolsep}{7pt}
\renewcommand{\arraystretch}{1.14}
\begin{tabular}{@{}l c c c c@{}}
\toprule
\rowcolor{headerblue}
\hd{Method} & \hd{Mean} & \hd{Worst frame} & \hd{10th pct.} &
\hd{$\mathrm{mean}-\mathrm{sd}$} \\
\midrule
\rowcolor{rowlight} \textbf{LIQE} & \pfB{11/15} & \pfB{11/15} & \pfB{11/15} & \pfB{11/15} \\
\rowcolor{rowwhite} NIQE          & \pfC{10/15} & \pfC{10/15} & \pfC{10/15} & \pfC{10/15} \\
\rowcolor{rowlight} QualiCLIP     & \pfD{\ 9/15} & \pfD{\ 9/15} & \pfB{11/15} & \pfC{10/15} \\
\rowcolor{rowwhite} BRISQUE       & \pfC{10/15} & \pfD{\ 9/15} & \pfD{\ 9/15} & \pfC{10/15} \\
\bottomrule
\end{tabular}
\end{table}


%% --- S10: tab:threshold_sens ---
\begin{table}[!ht]
\centering
\caption{Eligible set as a function of the generalization threshold, under both
readings of ``a database it was not fitted on''. \emph{Universal} requires the
threshold to hold on every unseen database (the reading adopted in the main article);
\emph{existential} requires it on at least one. The sensitivity requirement is
held at a within-reference mean above $0.50$ and at least $15/25$ types; the
topic-block lower confidence limit must also exceed the in-domain floor.}
\label{tab:threshold_sens}
\footnotesize
\setlength{\tabcolsep}{3pt}
\renewcommand{\arraystretch}{1.16}
\begin{tabular}{@{}c p{4.2cm} p{6.7cm}@{}}
\toprule
\rowcolor{headerblue}
\hd{Minimum $\rho^{*}$} & \hd{Universal reading (adopted)} & \hd{Existential reading} \\
\midrule
\rowcolor{rowlight} $0.50$ & LIQE, QualiCLIP, CLIP-IQA$+$, MUSIQ, TOPIQ, CLIP-IQA & LIQE, QualiCLIP, CLIP-IQA$+$, MUSIQ, TOPIQ, CLIP-IQA, DBCNN, HyperIQA \\
\rowcolor{rowwhite} $0.55$ & LIQE, QualiCLIP, CLIP-IQA$+$, MUSIQ & \emph{same eight methods} \\
\rowcolor{rowlight} $0.60$ & \textbf{LIQE, QualiCLIP, CLIP-IQA$+$} & \emph{same eight methods} \\
\rowcolor{perf5!16} $0.65^{\ast}$ & \textbf{LIQE, QualiCLIP, CLIP-IQA$+$} & \emph{same eight methods} \\
\rowcolor{rowwhite} $0.70$ & LIQE & \emph{same eight methods} \\
\rowcolor{rowlight} $0.75$ & \textcolor{accentred}{none} & LIQE, QualiCLIP, CLIP-IQA$+$, MUSIQ, TOPIQ, DBCNN, HyperIQA \\
\bottomrule
\multicolumn{3}{@{}p{0.95\linewidth}}{\scriptsize $^{\ast}$The value used.
At the adopted generalization threshold and universal reading, the same three
methods survive for sensitivity type-count cuts from $12/25$ through $22/25$;
none survives at $23/25$. Requiring only a positive floor point estimate instead
of a topic-block lower limit does not change the three-method set.} \\
\end{tabular}
\end{table}


%% --- S11: tab:e2e ---
\begin{table}[!ht]
\centering
\caption{The same comparison on the 15 human-annotated clusters. The two
references agree closely here; their disagreement emerges on the 42-cluster
analysis in the main article, which is the primary estimate. The DOVER column
uses the uncalibrated raw technical-plus-aesthetic sum, not DOVER-T.}
\label{tab:e2e}
\small
\setlength{\tabcolsep}{4.5pt}
\renewcommand{\arraystretch}{1.14}
\begin{tabular}{@{}l c r r r r@{}}
\toprule
\rowcolor{headerblue}
\hd{Scoring method} & \hd{Changed} & \hd{by raw T+A} & \hd{by Q-Align} &
\hd{Dur.\ (s)} & \hd{Words} \\
\midrule
\rowcolor{catbovw!12} SIFT+BoVW (deployed) & --- & \pfD{$52.1\%$} & \pfC{$61.1\%$} & $634.0$ & $1{,}613$ \\
\midrule
\rowcolor{perf5!16} \textbf{LIQE}  & $5/15$ & \pfA{$91.8\%$} & \pfA{$85.8\%$} & $455.6$ & $1{,}124$ \\
\rowcolor{rowwhite} BRISQUE        & $6/15$ & \pfA{$84.5\%$} & \pfA{$84.7\%$} & $431.0$ & $1{,}073$ \\
\rowcolor{rowlight} QualiCLIP      & $4/15$ & \pfA{$80.9\%$} & \pfA{$80.4\%$} & $439.6$ & $1{,}124$ \\
\rowcolor{rowwhite} NIQE           & $4/15$ & \pfA{$80.9\%$} & \pfA{$80.4\%$} & $439.6$ & $1{,}124$ \\
\midrule
\rowcolor{rowlight} DINOv2         & $12/15$ & \pfD{$53.4\%$} & \pfE{$36.7\%$} & $328.6$ & $\ \,904$ \\
\bottomrule
\end{tabular}
\end{table}

\FloatBarrier
\begin{center}
\textbf{Table S11 (continued). Construct and scale sensitivity on the 42
operational clusters.}\par\smallskip
\scriptsize
\setlength{\tabcolsep}{2.0pt}
\begin{tabular}{@{}L{1.55cm}L{2.25cm}rrL{2.6cm}L{2.25cm}@{}}
\toprule
\rowcolor{headerblue}
\hd{DOVER output} & \hd{Estimand} & \hd{SIFT} & \hd{LIQE} &
\hd{$\Delta$; topic CI} & \hd{Hierarchical CI} \\
\midrule
Technical branch & Global range (\%) & 50.5 & 76.0 & $+25.5$; $[-12.1,54.9]$ & $[-17.2,59.6]$ \\
\rowcolor{rowlight}Technical branch & Raw mean & $-.0770$ & $-.0701$ & $+.0069$; $[-.0028,.0178]$ & $[-.0041,.0197]$ \\
Technical branch & Mean cluster position (\%) & 57.1 & 64.4 & $+7.2$; $[-20.7,35.1]$ & $[-22.4,37.5]$ \\
\rowcolor{rowlight}Technical branch & Selected maximum (\%) & 54.8 & 59.5 & $+4.8$; $[-22.2,31.7]$ & $[-24.4,34.2]$ \\
Raw T+A sum & Global range (\%) & 47.4 & 77.7 & $+30.3$; $[1.1,54.1]$ & $[-4.1,59.0]$ \\
\rowcolor{rowlight}Aesthetic branch & Global range (\%) & 45.0 & 73.7 & $+28.7$; $[5.3,49.7]$ & $[-.2,54.6]$ \\
\bottomrule
\multicolumn{6}{@{}p{11.15cm}}{\scriptsize The branch sum is the raw
technical output plus the raw aesthetic output. It is not the official
ITU-standardized DOVER fused score. The technical branch is primary; the other
rows diagnose construct sensitivity.} \\
\end{tabular}

\medskip
\begin{tabular}{@{}lrrr@{}}
\toprule
\rowcolor{headerblue}\hd{DOVER-T range filter} & \hd{Clusters} &
\hd{Threshold} & \hd{Global-range $\Delta$} \\
\midrule
All clusters & 42 & --- & $+25.5$ \\
\rowcolor{rowlight}Exclude bottom 10\% & 37 & $.00397$ & $+25.9$ \\
Exclude bottom 25\% & 31 & $.00891$ & $+28.1$ \\
\bottomrule
\multicolumn{4}{@{}p{0.72\linewidth}}{\scriptsize DOVER-T cluster ranges:
minimum $.00028$, first quartile $.00891$, median $.01686$, third quartile
$.04118$, maximum $.10653$; zero-range clusters: 0.} \\
\end{tabular}

\medskip
\textbf{Matched-input control: ten identical frames and mean pooling.}\par\smallskip
\begin{tabular}{@{}lcrrr@{}}
\toprule
\rowcolor{headerblue}
\hd{Scorer} & \hd{Differs from SIFT} & \hd{DOVER-T range} & \hd{Raw mean} &
\hd{Selected max} \\
\midrule
SIFT+BoVW & --- & 52.0\% & $-.0766$ & 57.1\% \\
\rowcolor{rowlight}LIQE & 16/42 & 79.1\% & $-.0692$ & 64.3\% \\
\bottomrule
\multicolumn{5}{@{}p{0.93\linewidth}}{\scriptsize LIQE--SIFT global-range
$\Delta=+27.1$ points (topic CI $[5.6,47.8]$; hierarchical CI $[2.4,52.5]$).
Raw-branch $\Delta=+.00736$ (topic $[.00135,.01406]$; hierarchical
$[.00057,.01683]$). Thirteen of the 16 disagreement clusters belong to the
22-cluster complete-operation change set. Relative to the published deployed
path, matched-input SIFT changes 12 choices and matched-input LIQE changes 22;
the control isolates scorer identity, not the historical operation itself.
Matched SIFT uses a seeded per-cluster $K_{\mathrm{eff}}\leq300$ vocabulary,
segment-level document frequency and deployed IDF; empty frames remain in the
mean. Every segment yielded descriptors and no top-score ties occurred.} \\
\end{tabular}
\end{center}

%% --- S12: predictions recorded before the experiment ---
\begin{table}[!ht]
\centering
\caption{Eight predictions recorded in commit \texttt{9377c07}, whose Git
metadata predates the experiment, against measurements. The bundle verifies
content and metadata, not a third-party timestamp; none filters a method.}
\label{tab:preds_supp}
\footnotesize
\setlength{\tabcolsep}{4pt}
\begin{tabular}{@{}llp{4.6cm}p{3.4cm}l@{}}
\toprule
& \textbf{Subject} & \textbf{Prediction} & \textbf{Measured} & \textbf{Outcome} \\
\midrule
P1 & SIFT$+$BoVW & High accuracy on blur and darkening & $90\%$, $100\%$ & holds \\
P2 & SIFT$+$BoVW & Inverted on white noise & $8\%$ & holds \\
P3 & SIFT$+$BoVW & Near chance on JPEG and H.264 & $6\%$ and $90\%$ & refuted \\
P4 & CNN$+$BoVW ($19$) & Chance on every type & $44.9\%$, CI $[41.0,48.9]$; range $30$--$64$ & refuted \\
P5 & NIQE, BRISQUE & Accuracy falls under heavy overlay & $+2.0$, $+2.2$ points & refuted \\
P6 & Learned methods & No interaction with overlay & $-18.1$ to $-8.1$ for three & refuted \\
P7 & SIFT$+$BoVW & Design~B mean below $0.15$ & $-0.425$ & refuted \\
P8 & Q-Align & Design~A accuracy above $0.90$ & $98.2\%$, CI $[97.2,99.1]$ & holds \\
\bottomrule
\end{tabular}
\end{table}

The complete prediction-recording commit is:
\par\smallskip\noindent\path{9377c07ca080ac61c7cff80e0ff6c7507b6ebaba}.
\par\smallskip The archive's \texttt{registered\_prediction\_commit.bundle}
makes the object and metadata inspectable; Git dates are user-supplied fields,
not an externally certified timestamp.

\FloatBarrier
\begin{center}
\textbf{Table S12 (continued). Fidelity and descriptor-count diagnostics for the
constructed broadcast pairs.}\par\smallskip
\scriptsize
\setlength{\tabcolsep}{5pt}
\begin{tabular}{@{}lrrrrrr@{}}
\toprule
\rowcolor{headerblue}
\hd{Method} & \multicolumn{3}{c}{\cellcolor{headerblue}\hd{All 8 types (800)}} &
\multicolumn{3}{c}{\cellcolor{headerblue}\hd{Without dark/desaturation (600)}} \\
\rowcolor{headerblue}
& \hd{Correct} & \hd{Ties} & \hd{Accuracy} & \hd{Correct} & \hd{Ties} & \hd{Accuracy} \\
\midrule
SIFT+BoVW & 521 & 0 & 65.1\% & 382 & 0 & 63.7\% \\
\rowcolor{rowlight}LIQE & 796 & 0 & 99.5\% & 599 & 0 & 99.8\% \\
Q-Align & 786 & 0 & 98.2\% & 600 & 0 & 100.0\% \\
\rowcolor{rowlight}DOVER-T & 781 & 0 & 97.6\% & 599 & 0 & 99.8\% \\
\bottomrule
\end{tabular}

\medskip
\begin{tabular}{@{}lrrp{5.6cm}@{}}
\toprule
\rowcolor{headerblue}\hd{Type} & \hd{Mean $\Delta D$} & \hd{$\Delta D>0$} &
\hd{Mean $\Delta D$ by level 1--5} \\
\midrule
Blur & $-797$ & 10\% & $-130,-417,-913,-1188,-1339$ \\
\rowcolor{rowlight}Dark & $-796$ & 0\% & $-229,-487,-791,-1116,-1358$ \\
H.264 & $-398$ & 10\% & $-128,-220,-361,-541,-737$ \\
\rowcolor{rowlight}JPEG & $+289$ & 95\% & $+161,+252,+339,+376,+317$ \\
Motion blur & $-1067$ & 0\% & $-548,-909,-1174,-1308,-1396$ \\
\rowcolor{rowlight}Noise & $+469$ & 92\% & $+126,+348,+517,+607,+749$ \\
Rescale & $-646$ & 12\% & $-154,-377,-585,-911,-1204$ \\
\rowcolor{rowlight}Desaturation & $+18$ & 61\% & $0,+14,+21,+34,+24$ \\
\bottomrule
\multicolumn{4}{@{}p{0.9\linewidth}}{\scriptsize $\Delta D=D_{\mathrm{degraded}}-
D_{\mathrm{unmodified}}$. Strict accuracy uses correct/total; a tie is neither
correct nor incorrect and remains in the denominator.} \\
\end{tabular}

\medskip
\setlength{\tabcolsep}{2pt}
\begin{tabular}{@{}lrrrrrrrrr@{}}
\toprule
\rowcolor{headerblue}
\hd{DOVER-T} & \hd{dark} & \hd{motion} & \hd{blur} & \hd{H.264} &
\hd{rescale} & \hd{satur.} & \hd{noise} & \hd{JPEG} & \hd{All} \\
\midrule
Accuracy (\%) & 93 & 100 & 100 & 100 & 100 & 89 & 99 & 100 & 97.6 \\
\rowcolor{rowlight}Mean severity $\rho^{*}$ & .985 & .900 & 1.000 & 1.000 &
1.000 & .715 & .995 & 1.000 & --- \\
\bottomrule
\multicolumn{10}{@{}p{0.95\linewidth}}{\scriptsize The raw technical output is
the exact branch later used for the 42-cluster output analysis. Its overall
segment-block 95\% CI is $[96.1,99.0]$; after excluding darkening and
desaturation, accuracy is 99.8\% with CI $[99.5,100]$. There are no ties. Each
condition is encoded through the same 5-fps MP4 wrapper and uses paired spatial
fragment seed 42.} \\
\end{tabular}
\end{center}

%% --- S13: method inventory ---
\begin{table}[!ht]
\centering
\caption{Method inventory and released-weight supervision. \emph{Sup.}:
OU = opinion-unaware, K = KonIQ-10k, L = LIVE, M = multi-database mixture.
Roles follow Table~1: D = deployed, C = candidate, R = reference and X =
diagnostic or ablation. The daggered partial-coverage methods are excluded from
complete-case screening.}
\label{tab:methods_supp}
\scriptsize
\setlength{\tabcolsep}{2.5pt}
\begin{tabular}{@{}p{2.0cm}p{2.0cm}p{3.0cm}cccc@{}}
\toprule
\rowcolor{headerblue}
\hd{Method} & \hd{Family} & \hd{Construct} & \hd{Sup.} & \hd{Pol.} & \hd{GPU} & \hd{Role} \\
\midrule
SIFT+BoVW & Handcrafted & TF-IDF word count & -- & \polup & No & D \\
Keypoint density & Handcrafted & keypoints/frame & -- & \polup & No & X \\
DINOv2 & Self-supervised & patch-token norm & -- & \polup & Yes & X \\
ResNet-50 & Supervised & pooled-feature norm & -- & \polup & Yes & X \\
BRISQUE & NSS & natural-statistics deviation & L & \poldown & No & C \\
NIQE & NSS & pristine-model distance & OU & \poldown & No & C \\
QualiCLIP & CLIP & quality score & OU & \polup & Yes & C \\
LIQE & CLIP & quality, distortion, scene & K & \polup & Yes & C \\
Q-Align & LMM & textual quality levels & M & \polup & Yes & R \\
SI & ITU-T complexity & Sobel energy & -- & \polup & No & X \\
CNN+BoVW (19 extractors) & ablation & TF-IDF word count & -- & \polup & Yes & X \\
DBCNN & CNN regression & bilinear pooling & K & \polup & Yes & C \\
HyperIQA & CNN regression & adaptive weights & K & \polup & Yes & C \\
MUSIQ & Transformer & multi-scale score & K & \polup & Yes & C \\
TOPIQ & Transformer & semantic attention & K & \polup & Yes & C \\
CLIP-IQA & CLIP & antonym prompts & OU & \polup & Yes & C \\
CLIP-IQA$+$ & CLIP & tuned prompts & K & \polup & Yes & C \\
ARNIQA$^{\dagger}$ & Self-supervised & distortion manifold & K & \polup & Yes & partial \\
TReS$^{\dagger}$ & CNN--transformer & ranking/consistency & K & \polup & Yes & partial \\
\midrule
\multicolumn{7}{@{}p{11.2cm}}{\scriptsize PyIQA registry keys used by the
archived scorers are \texttt{brisque}, \texttt{niqe}, \texttt{dbcnn},
\texttt{hyperiqa}, \texttt{musiq}, \texttt{topiq\_nr}, \texttt{clipiqa},
\texttt{clipiqa+}, \texttt{liqe}, \texttt{qualiclip}, \texttt{arniqa} and
\texttt{tres}. The final inventory and analysis archive were frozen on
23 August 2026. MANIQA was attempted with PyIQA 0.1.15.post2 and the
\texttt{maniqa-kadid} checkpoint. A single-image diagnostic completed, but the
released implementation evaluates 20 crops per image and the full 21,360-image
three-database run did not finish within the experiment window on the single
8-GB RTX~4060; it is therefore an uncompleted method, not a failed candidate.} \\
\bottomrule
\end{tabular}
\end{table}

\addtocounter{table}{-1}
\renewcommand{\theHtable}{S13-extractors}
\begin{table}[!ht]
\centering
\caption{(continued). Complete extractor inventory for the
CNN+BoVW ablation. Every row uses the final spatial tensor before global
pooling or classification.}
\label{tab:cnn_inventory_supp}
\scriptsize
\setlength{\tabcolsep}{3.2pt}
\renewcommand{\arraystretch}{1.04}
\begin{tabular}{@{}p{2.65cm}p{3.45cm}p{4.35cm}@{}}
\toprule
\rowcolor{headerblue}
\hd{Extractor} & \hd{Released torchvision weights} & \hd{Spatial endpoint} \\
\midrule
ResNet-18 & IMAGENET1K\_V1 & trunk before average pooling \\
\rowcolor{rowlight}ResNet-50 & IMAGENET1K\_V2 & trunk before average pooling \\
ResNet-101 & IMAGENET1K\_V2 & trunk before average pooling \\
\rowcolor{rowlight}VGG-16 & IMAGENET1K\_V1 & full \texttt{features} module \\
VGG-19 & IMAGENET1K\_V1 & full \texttt{features} module \\
\rowcolor{rowlight}EfficientNet-B0 & IMAGENET1K\_V1 & full \texttt{features} module \\
MobileNet-V3-Large & IMAGENET1K\_V2 & full \texttt{features} module \\
\rowcolor{rowlight}ConvNeXt-Tiny & IMAGENET1K\_V1 & full \texttt{features} module \\
ConvNeXt-Small & IMAGENET1K\_V1 & full \texttt{features} module \\
\rowcolor{rowlight}DenseNet-121 & IMAGENET1K\_V1 & full \texttt{features} module \\
DenseNet-201 & IMAGENET1K\_V1 & full \texttt{features} module \\
\rowcolor{rowlight}ResNeXt-50-32x4d & IMAGENET1K\_V2 & trunk before average pooling \\
ResNeXt-101-64x4d & IMAGENET1K\_V1 & trunk before average pooling \\
\rowcolor{rowlight}Wide-ResNet-50-2 & IMAGENET1K\_V2 & trunk before average pooling \\
Wide-ResNet-101-2 & IMAGENET1K\_V2 & trunk before average pooling \\
\rowcolor{rowlight}Inception-V3 & IMAGENET1K\_V1 & convolutional trunk through Mixed\_7c \\
EfficientNet-B4 & IMAGENET1K\_V1 & full \texttt{features} module \\
\rowcolor{rowlight}SqueezeNet-1.1 & IMAGENET1K\_V1 & full \texttt{features} module \\
MobileNet-V2 & IMAGENET1K\_V2 & full \texttt{features} module \\
\bottomrule
\multicolumn{3}{@{}p{0.96\linewidth}}{\scriptsize Common preprocessing:
bicubic resize to 256 pixels, centre crop to $224\times224$, ImageNet mean and
standard deviation, and FP16 inference. Each spatial location is flattened into
one descriptor and $\ell_2$-normalised in FP32. The same internal key-frame
filter, vocabulary construction, TF--IDF weighting and segment aggregation are
then applied. The archived implementation is \texttt{cnn\_bovw\_scorer.py}.} \\
\end{tabular}
\end{table}

\addtocounter{table}{-1}
\renewcommand{\theHtable}{S13-references}
\begin{table}[!ht]
\centering
\caption{(continued). Exact configurations of the two model-based
references used in the broadcast analyses.}
\label{tab:reference_config_supp}
\scriptsize
\setlength{\tabcolsep}{2.5pt}
\begin{tabular}{@{}L{1.2cm}L{2.05cm}L{3.0cm}L{4.05cm}@{}}
\toprule
\rowcolor{headerblue}
\hd{Model} & \hd{Config.} & \hd{Input / pooling} & \hd{Output used} \\
\midrule
Q-Align & \texttt{q-future/}\allowbreak\texttt{one-align}; task \texttt{quality}, input
\texttt{image} & Uniform 1-fps frames; arithmetic mean of frame scores &
Four-bit NF4 with double quantisation and FP16 compute; higher score is better.
The same configuration scores the 800 constructed pairs and supplies the
in-domain floor. \\
\rowcolor{rowlight}DOVER & Official \texttt{DOVER.pth} and
\texttt{dover.yml}, \texttt{val-}\allowbreak\texttt{l1080p} view decomposition & Technical:
$7\times7$ spatial fragments of size 32, 32 frames, interval 2, three clips;
aesthetic: $224\times224$, 32 frames, interval 2, one clip & The technical
branch is the primary output reference and is also evaluated on the 800 pairs
with a common 5-fps MP4 wrapper and paired seed 42. The raw branch sum is only a
sensitivity analysis; it is not the official standardised fusion or DOVER++.
Higher (less negative) raw technical output indicates better quality. \\
\bottomrule
\end{tabular}
\end{table}
\renewcommand{\theHtable}{S\arabic{table}}

%% --- S14: computational cost ---
\begin{table}[!ht]
\centering
\caption{Computational cost measured on one RTX~4060 over the same inputs.
Frame-wise entries report median latency after ten warm-up and sixty timed
iterations; the final two rows are measured complete paths.}
\label{tab:cost_supp}
\small
\setlength{\tabcolsep}{6pt}
\begin{tabular}{@{}lrrl@{}}
\toprule
\rowcolor{headerblue}
\hd{Method} & \hd{ms/frame} & \hd{IQR} & \hd{Per cluster} \\
\midrule
SI & 2.2 & 0.14 & 0.2~s \\
QualiCLIP & 14.3 & 1.68 & 1.2~s \\
TOPIQ & 18.7 & 0.36 & 1.5~s \\
BRISQUE & 19.6 & 0.98 & 1.6~s \\
CLIP-IQA$+$ & 21.4 & 0.65 & 1.8~s \\
DBCNN & 21.5 & 0.46 & 1.8~s \\
LIQE & 22.7 & 0.82 & 1.9~s \\
CLIP-IQA & 25.2 & 0.51 & 2.1~s \\
NIQE & 34.6 & 0.90 & 2.8~s \\
MUSIQ & 40.0 & 9.36 & 3.3~s \\
ResNet-50 & 46.2 & 8.44 & 3.8~s \\
DINOv2 & 46.4 & 36.63 & 3.8~s \\
Keypoint density & 46.8 & 1.48 & 3.8~s \\
HyperIQA & 50.1 & 0.77 & 4.1~s \\
\midrule
SIFT+BoVW & $\sim1{,}724$ & --- & 182~s \\
Q-Align & 322.6 & --- & 26~s \\
\bottomrule
\multicolumn{4}{@{}p{0.94\linewidth}}{\scriptsize The uniform 1-fps candidate
cache contains 2,892 frames from 83 distinct segments. Counting repeated
memberships across the 93 cluster slots gives 3,525 segment-frame instances.
Candidate latency was timed per decoded frame; the SIFT+BoVW row is the complete
operation divided by its input-frame count and therefore includes key-frame
selection, descriptor extraction, vocabulary quantisation and aggregation. The
per-cluster column is the operational quantity; the per-frame SIFT value is an
approximate normalisation rather than an identical-path microbenchmark.} \\
\end{tabular}
\end{table}

%% --- S15: descriptive training-database and database-regime contrasts ---
\begin{table}[!ht]
\centering
\caption{Descriptive three-database contrasts and released-weight exposure.
K--C is KonIQ-10k minus CLIVE; C--KADID is CLIVE minus KADID-10k. Positive
values are losses. K = KonIQ-10k opinion training, OU = opinion-unaware and L =
LIVE opinion training. The contrasts do not isolate causal effects because the
databases, architectures and objectives are unmatched.}
\label{tab:exposure_supp}
\footnotesize
\setlength{\tabcolsep}{4.5pt}
\renewcommand{\arraystretch}{0.98}
\begin{tabular}{@{}lcrrrrr@{}}
\toprule
\rowcolor{headerblue}
\hd{Method} & \hd{Sup.} & \hd{KonIQ} & \hd{CLIVE} & \hd{KADID} &
\hd{K--C} & \hd{C--KADID} \\
\midrule
\multicolumn{7}{@{}l}{\itshape\textcolor{headerblue}{Released checkpoint fitted on KonIQ-10k}} \\
\rowcolor{rowlight} TOPIQ & K & 0.973$^{\dagger}$ & 0.811 & 0.511 & $+0.162$ & $+0.300$ \\
\rowcolor{rowwhite} HyperIQA & K & 0.947$^{\dagger}$ & 0.750 & 0.462 & $+0.197$ & $+0.288$ \\
\rowcolor{rowlight} DBCNN & K & 0.919$^{\dagger}$ & 0.756 & 0.484 & $+0.163$ & $+0.272$ \\
\rowcolor{rowwhite} MUSIQ & K & 0.873$^{\dagger}$ & 0.789 & 0.556 & $+0.084$ & $+0.233$ \\
\rowcolor{rowlight} CLIP-IQA$+$ & K & 0.824$^{\dagger}$ & 0.805 & 0.654 & $+0.019$ & $+0.151$ \\
\rowcolor{rowwhite} LIQE & K & 0.915$^{\dagger}$ & 0.865 & 0.723 & $+0.050$ & $+0.142$ \\
\rowcolor{gray!10} \multicolumn{5}{@{}l}{\itshape mean over the six} & $+0.113$ & $+0.231$ \\
\midrule
\multicolumn{7}{@{}l}{\itshape\textcolor{headerblue}{Never fitted on KonIQ-10k}} \\
\rowcolor{rowlight} QualiCLIP & OU & 0.711 & 0.757 & 0.665 & $-0.046$ & $+0.092$ \\
\rowcolor{rowwhite} CLIP-IQA & OU & 0.647 & 0.703 & 0.501 & $-0.056$ & $+0.202$ \\
\rowcolor{rowlight} NIQE & OU & 0.368 & 0.449 & 0.379 & $-0.081$ & $+0.070$ \\
\rowcolor{rowwhite} BRISQUE & L & 0.236 & 0.313 & 0.330 & $-0.077$ & $-0.017$ \\
\rowcolor{gray!10} \multicolumn{5}{@{}l}{\itshape mean over the four} & $-0.065$ & $+0.087$ \\
\bottomrule
\multicolumn{7}{@{}p{0.95\linewidth}}{\scriptsize $^{\dagger}$In-distribution.
CLIVE contains 1,162 images and is unseen by every checkpoint in the table. The
operational partition assigns 33\% of the mean KonIQ-to-KADID gap of the fitted
group to K--C and 67\% to C--KADID; these percentages describe contrasts and are
not causal estimates.} \\
\end{tabular}
\end{table}

%% --- S16: detectable-keypoint conditioning ---
\begin{table}[!ht]
\centering
\caption{Spearman association with human scores before and after conditioning
on detectable keypoint density. $r_{\mathrm{kp}}$ is the rank correlation
between the method score and keypoint density. Differences are descriptive
attenuations, not causal fractions of ``content density.''}
\label{tab:partial_supp}
\footnotesize
\setlength{\tabcolsep}{2.5pt}
\renewcommand{\arraystretch}{0.98}
\begin{tabular}{@{}lrrrrrrrr@{}}
\toprule
\rowcolor{headerblue}
\hd{Method} & \multicolumn{4}{c}{\cellcolor{headerblue}\hd{KonIQ-10k}} &
\multicolumn{4}{c}{\cellcolor{headerblue}\hd{KADID-10k}} \\
\rowcolor{headerblue}
& \hd{$r_{\mathrm{kp}}$} & \hd{$|\rho|$} & \hd{$|\rho_{\cdot\mathrm{kp}}|$} & \hd{Diff.}
& \hd{$r_{\mathrm{kp}}$} & \hd{$|\rho|$} & \hd{$|\rho_{\cdot\mathrm{kp}}|$} & \hd{Diff.} \\
\midrule
Q-Align & $+.44$ & .920 & .901 & $-.019$ & $+.13$ & .951 & .950 & $-.001$ \\
LIQE & $+.43$ & .915 & .894 & $-.020$ & $+.18$ & .723 & .718 & $-.005$ \\
TOPIQ & $+.48$ & .973 & .966 & $-.007$ & $+.21$ & .511 & .500 & $-.011$ \\
HyperIQA & $+.48$ & .947 & .933 & $-.014$ & $+.17$ & .462 & .451 & $-.011$ \\
DBCNN & $+.49$ & .919 & .898 & $-.021$ & $+.22$ & .484 & .472 & $-.012$ \\
MUSIQ & $+.44$ & .873 & .841 & $-.032$ & $+.20$ & .556 & .546 & $-.010$ \\
QualiCLIP & $+.53$ & .711 & .626 & $-.086$ & $+.29$ & .665 & .663 & $-.002$ \\
CLIP-IQA & $+.27$ & .647 & .611 & $-.036$ & $+.22$ & .501 & .489 & $-.012$ \\
\midrule
SI & $+.77$ & .513 & .296 & $-.218$ & $+.70$ & .219 & .187 & $-.032$ \\
NIQE & $-.63$ & .368 & .125 & $-.243$ & $-.33$ & .379 & .362 & $-.018$ \\
BRISQUE & $-.31$ & .236 & .116 & $-.120$ & $-.13$ & .330 & .319 & $-.011$ \\
\midrule
\rowcolor{gray!14} SIFT+BoVW (canonical) & $+.998$ & .447 & .016 & arith. & $+.992$ & .125 & .028 & arith. \\
\rowcolor{gray!14} SIFT+BoVW (deployed) & $-.9999$ & .445 & n.d. & arith. & $-.9998$ & .123 & n.d. & arith. \\
\bottomrule
\multicolumn{9}{@{}p{0.96\linewidth}}{\scriptsize The shaded rows are the
near-tautological regime: conditioning a score on a variable with which it is
almost perfectly correlated removes the variable from itself. For the deployed
row the partial coefficient is not reported because the denominator approaches
zero and the result is numerically unstable.} \\
\end{tabular}
\end{table}

\end{document}
